\chapter{Fazit und Ausblick}
\section{Fazit}
Ziel dieser Bachelorarbeit war die Einführung eines Standards zur Aufbereitung und Auswertung von Messdaten der Hydropulsanlage. Hierzu wurde eine MATLAB-basierte Auswertesoftware entwickelt, mit der die bei Schwingfestigkeitsversuchen entstehenden Rohmessdaten weitgehend automatisiert verarbeitet und entsprechend der DIN~50100 ausgewertet werden können. Neben der eigentlichen Versuchsauswertung standen insbesondere eine Reduzierung des manuellen Auswerteaufwands, eine reproduzierbare Vorgehensweise sowie eine möglichst bedienfreundliche Anwendung im Vordergrund.

Im Rahmen der Arbeit wurde dafür ein modular aufgebauter Programmablauf entwickelt, der die einzelnen Schritte von der Eingabe der Versuchsparameter über das Einlesen und die Aufbereitung der Messdaten bis hin zur statistischen Auswertung miteinander verbindet. Die Messdateien des Zeit- und Langzeitfestigkeitsbereichs werden getrennt eingelesen und anschließend automatisiert weiterverarbeitet. Dabei werden die für die Auswertung benötigten Messgrößen Kraft, Weg, Zeit und Schwingspielzahl unabhängig von ihrer Position innerhalb der Messdateien identifiziert. Durch die Berücksichtigung der unterschiedlichen Kraft- und Wegsensoren der drei Hydropulsanlagen ist die Auswertung nicht auf eine einzelne Prüfmaschine beschränkt.

Auf Grundlage der aufbereiteten Messdaten erfolgt anschließend die Auswertung des 
Zeit- und Langzeitfestigkeitsbereichs. Für die Zeitfestigkeit wird das 
Perlenschnurverfahren mit einer linearen Regression der logarithmierten Versuchswerte 
angewendet, während die Langzeitfestigkeit mithilfe des Treppenstufenverfahrens 
bestimmt wird. Neben der 50\,\%-Wöhlerlinie werden die Streuungen der Versuchsergebnisse 
durch die 10\,\%- und 90\,\%-Wöhlerlinien berücksichtigt. Über die Bestimmung des 
Knickpunktes werden der Zeit- und Langzeitfestigkeitsbereich schließlich zu einer 
vollständigen Wöhlerlinie zusammengeführt.

Ein weiterer Schwerpunkt lag auf der Automatisierung und Standardisierung des gesamten 
Auswerteprozesses. Plausibilitätsprüfungen und Fehlermeldungen sollen fehlerhafte 
Eingaben oder für die Auswertung ungeeignete Datensätze frühzeitig erkennen. Dadurch 
wird nicht nur der manuelle Aufwand reduziert, sondern auch die Nachvollziehbarkeit 
und Reproduzierbarkeit der Versuchsauswertung verbessert. Die Funktionsfähigkeit der 
implementierten Berechnungen wurde anhand von realen sowie realitätsnah generierten 
Messdatensätzen überprüft. Zusätzlich wurden ausgewählte Berechnungsschritte 
exemplarisch nachvollzogen, wobei sich gegenüber den MATLAB-Ergebnissen lediglich 
geringe, auf Rundungen zurückzuführende Abweichungen ergaben.

Abschließend werden die ermittelten Ergebnisse automatisiert in Form der Wöhlerlinie 
dargestellt und in einem Prüfprotokoll zusammengefasst. Damit steht ein durchgängiger 
Auswerteprozess zur Verfügung, der die Verarbeitung der Rohmessdaten, deren Auswertung, 
die Visualisierung der Ergebnisse und deren Dokumentation miteinander verbindet. 
Der entwickelte modulare Aufbau ermöglicht darüber hinaus eine zukünftige Anpassung 
und Erweiterung der Auswertesoftware. Insgesamt konnte somit ein standardisierter 
und weitgehend automatisierter Ablauf für die Auswertung der Schwingfestigkeitsversuche 
an den Hydropulsanlagen geschaffen werden.

\section{Ausblick}

Die im Rahmen dieser Arbeit entwickelte Auswertesoftware bildet eine Grundlage für 
die standardisierte Auswertung zukünftiger Schwingfestigkeitsversuche an den 
Hydropulsanlagen. Für die weitere Entwicklung bietet sich insbesondere eine 
Validierung der Software anhand zusätzlicher realer Versuchsreihen an. Dadurch 
können die implementierten Auswerteverfahren für unterschiedliche Werkstoffe, 
Probengeometrien und Versuchsparameter überprüft und die Anwendungsgrenzen der 
Software weiter untersucht werden.

Darüber hinaus kann die Auswertesoftware zukünftig um weitere statistische 
Auswerteverfahren ergänzt werden. Durch den modularen und funktionsbasierten Aufbau 
des Programmcodes können zusätzliche Berechnungs- und Auswerteschritte in den 
bestehenden Programmablauf integriert werden, ohne die grundlegende Struktur der 
Software verändern zu müssen. Ebenso können weitere Plausibilitätsprüfungen 
implementiert werden, um fehlerhafte oder ungewöhnliche Messdaten bereits während 
der automatisierten Verarbeitung zu erkennen.

Weiteres Entwicklungspotenzial besteht bei der automatisierten Erstellung des 
Prüfprotokolls. Neben einer Erweiterung der ausgegebenen Versuchsinformationen und 
Kennwerte könnten zukünftig zusätzliche Diagramme oder Auswertungen automatisch in 
das Prüfprotokoll integriert werden. Auf diese Weise kann der in dieser Arbeit 
entwickelte Standard schrittweise erweitert und an zukünftige Anforderungen bei 
der Durchführung und Auswertung von Schwingfestigkeitsversuchen angepasst werden.
